\section{Introduction}\label{sec:intro}

The standard practice in volatility-forecasting research is to compute a loss function on each of several candidate forecasters, rank them, and report the lowest-loss model as the winner. This practice has been the convention for at least three decades. It is also, in a meaningful fraction of published comparisons, statistically indefensible. The headline gap between the winner and runner-up is often smaller than the standard error of the loss differential; the ``win'' that appears decisive in a clean table does not survive a basic Diebold-Mariano test.

The methodological tools to do this correctly have existed since \citet{DieboldMariano1995}. Refinements that handle multiple comparisons \citep{HansenLundeNason2011}, that test a benchmark against a set of competitors \citep{Hansen2005}, and that protect against data-snooping bias when many forecasters are tried \citep{White2000} are all between fifteen and thirty years old. The tools are not new. What is missing in the practice of the volatility-forecasting subfield is the discipline of applying them by default.

We make two contributions. First, we introduce \texttt{vol-eval}, an open-source Python package that bundles the standard loss functions and significance tests behind a clean, functional API designed to make rigorous evaluation a one-line addition to any existing forecasting workflow. The package is MIT-licensed, has 34 unit tests, and is cross-validated for numerical correctness against the canonical \texttt{arch} package \citep{Sheppard2023}. Second, we demonstrate the package by applying it to the seven-model volatility-forecasting horse race of \citet{KhanVol2026}, a controlled comparison on 980 trading days of S\&P 500 5-day realized volatility. The demonstration produces three findings that we view as representative of what a fuller re-analysis of the published literature would reveal:

\begin{enumerate}[leftmargin=*]
    \item Of 21 ordered pairs of models, the Diebold-Mariano test rejects equal accuracy at the 5\% level for only 6 pairs. The headline ``ensemble narrowly wins overall'' result is not statistically distinguishable from the runner-up GJR-GARCH at $p = 0.90$.
    \item The 90\% Model Confidence Set keeps six of the seven models. Only HAR-RV is eliminated. The notion that any single forecaster is the ``winner'' of this seven-model comparison is not statistically supportable.
    \item The Hansen SPA test, run with each model in turn as the benchmark, rejects non-inferiority at the 5\% level only when HAR-RV is the benchmark. The other six models cannot be rejected as the best forecaster in the panel.
\end{enumerate}

These findings are produced by \texttt{vol-eval} in a few seconds of laptop compute. They were available, in principle, to any reader of \citet{KhanVol2026}, but they appeared only as a partial Diebold-Mariano table in the original paper's \S6.3 because no convenient packaged tooling existed for the broader analysis. The discipline of running the full suite of tests is the discipline that an off-the-shelf package makes affordable.

\subsection{Why publish a tool now}

A natural question for a methods paper is: why now? Three reasons.

\textit{The audience is in transition.} The capital base for ML-driven volatility forecasting is moving toward institutions with explicit model-risk obligations: banks under SR 11-7 \citep{FedSR117} and OCC 2011-12 guidance, EU institutions under the AI Act \citep{EUAIAct}, insurance companies under their own model-validation regimes. These audiences will increasingly demand audit-grade documentation of forecast comparisons, including statistical significance of any reported ``win.'' The free-form practices that worked when the audience was hedge-fund quants do not satisfy a model-validation analyst.

\textit{The publication norm is moving.} Editorial calls for more rigorous reporting in financial machine learning have become common \citep{ArnottHarveyMarkowitz2019, Lopez2018}. \citet{Khan2026Survey} documents that 84\% of empirical financial-ML papers from 2015--2025 do not ship code under the basic Reproducibility Disclosure Score rubric, and that the disclosure rate has \emph{fallen} since 2020 despite a decade of editorial pressure. Significance reporting is in roughly the same state: the modal paper reports a clean ranking with no test. The norm will move because regulators and large-institution adopters will eventually require it. The tooling needs to exist when that demand materializes.

\textit{The marginal cost of doing it right is now small.} Implementing the Diebold-Mariano test with a Newey-West standard error is about thirty lines of Python. The Model Confidence Set is a few hundred. The Hansen SPA with three p-value variants is comparable. None of this is hard. The barrier has been the cost of writing it from scratch each time, plus the temptation to skip it because reviewers do not require it. A maintained package removes the writing cost; the demonstration paper attached to that package makes the case for not skipping.

\subsection{What this paper is not}

This paper is a methods paper with a self-contained empirical demonstration on a single seven-model panel. It is not a re-analysis of the published volatility-forecasting literature. We make the case for such a re-analysis in \Cref{sec:discussion}, and the project repository contains the data-collection scaffolding for a 30-paper sample, but the empirical extension to that sample is queued as Phase 2 of the project and is the subject of a follow-on paper currently in preparation. Our headline findings here are about the seven-model comparison of \citet{KhanVol2026} and should not be over-generalized.

\subsection{Roadmap}

\Cref{sec:package} describes the \texttt{vol-eval} package: design choices, API, dependencies, test coverage. \Cref{sec:methodology} reviews the four significance tests we implement, with attention to the implementation details that matter for reproducibility (HAC bandwidth selection, bootstrap design, centering choices in SPA). \Cref{sec:demonstration} applies the package to the seven-model panel of \citet{KhanVol2026} and reports the three findings above in detail. \Cref{sec:discussion} interprets the demonstration and lays out the queued Phase 2 extension. \Cref{sec:reproducibility} states the reproducibility commitment.
